% main.tex -- Paper 2, IEEE Transactions on Neural Networks and Learning Systems (Full Paper), IEEEtran journal mode.
% Build skeleton created 2026-10-29. NO PAPER CONTENT: every section is an \input stub with TODO comments.
% Requirements and their sources: research/paper2/tnnls/TNNLS-CHECKLIST-2026-10-29.md  (read it before editing the layout).
%
% Switches (set near the top of this file):
%   \anontrue / \anonfalse   double-anonymous review build (default: anonymous) vs. camera-ready with authors, funding, biographies.
%                            The CIS page says double-anonymous, the Xplore page says reviewers know the authors (checklist C1).
%   \placeholdertrue / false \nocite{*} so the placeholder PDF lists the whole reference file (tests blue links, column balance).
%                            Set to false as soon as real \cite commands exist.
%
% Limits (CIS Information for Authors, section H): Full Paper 10 pages free, hard maximum 15 pages (USD 200 per page above 10),
% supplement excluded from the count. Abstract 150-250 words, no abbreviations or maths; 4-5 Index Terms.
\documentclass[lettersize,journal]{IEEEtran}   % IEEE Template Selector TNNLS LaTeX zip (2025-10-24): local IEEEtran.cls = IEEE copy of v1.8b with its 2023 title-size edits

% ---- switches -----------------------------------------------------------------------------------
\newif\ifanon      \anontrue
\newif\ifplaceholder \placeholdertrue

% ---- packages -----------------------------------------------------------------------------------
\usepackage[T1]{fontenc}
\usepackage{lmodern}                 % Type 1 fonts, embedded in the PDF (journal requires all fonts embedded)
\usepackage{amsmath,amssymb,amsthm}  % IEEE Mathematics Guide: numbered equations, consecutive numbering
\usepackage{mathtools}
\usepackage{graphicx}
\usepackage{booktabs}
\usepackage{xcolor}
\usepackage[ruled,vlined,linesnumbered]{algorithm2e}   % Input/Output/Data blocks; build up across the paper
\usepackage{cite}                    % IEEEtran's recommended citation grouping [1]-[3]
\usepackage{microtype}
\usepackage{url}
\def\UrlBreaks{\do\/\do-\do.\do=\do_\do?\do\&\do\%\do\#}   % let long DOI/URL strings break in a 3.5 in column
\usepackage[colorlinks=true,linkcolor=blue,citecolor=blue,urlcolor=blue,filecolor=blue,
            bookmarksnumbered=true,pdfusetitle,breaklinks=true]{hyperref}   % all links blue (Shlok, memory: paper-presentation-rules)
\usepackage[capitalise,nameinlink]{cleveref}   % after hyperref
\crefname{figure}{Fig.}{Figs.}          % IEEE: "Fig. 1", abbreviation also at sentence start
\Crefname{figure}{Fig.}{Figs.}
\crefname{table}{Table}{Tables}
\crefformat{equation}{(#2#1#3)}
\crefmultiformat{equation}{(#2#1#3)}{ and~(#2#1#3)}{, (#2#1#3)}{, and~(#2#1#3)}
\crefname{algocf}{Algorithm}{Algorithms}
\crefname{section}{Section}{Sections}
\crefname{appendix}{Appendix}{Appendices}

\interdisplaylinepenalty=2500        % IEEEtran recommendation for amsmath
\hyphenation{op-tical net-works semi-conduc-tor}
\emergencystretch=2em

% ---- theorem environments (amsthm). IEEEtran supplies IEEEproof; amsthm's proof is also fine. -----
\theoremstyle{plain}
\newtheorem{theorem}{Theorem}
\newtheorem{proposition}{Proposition}
\newtheorem{lemma}{Lemma}
\newtheorem{corollary}{Corollary}
\theoremstyle{definition}
\newtheorem{definition}{Definition}
\newtheorem{assumption}{Assumption}
\theoremstyle{remark}
\newtheorem{remark}{Remark}

% ---- project macros -----------------------------------------------------------------------------
% \todo marks an open fact: [AUTHOR: ...] = fact only the authors can supply; [RESULT: ...] = result that does not exist yet.
% All \todo items must be resolved before submission (the build prints a count).
\newcommand{\todo}[1]{{\color{red}\textbf{[#1]}}}
\newcommand{\E}{\mathbb{E}}
\newcommand{\PP}{\mathbb{P}}
\newcommand{\ind}{\mathbb{1}}
\newcommand{\Var}{\operatorname{Var}}
% anonymised link: shown as plain text in the anonymous build, a real link in the camera-ready build
\newcommand{\anonlink}[2]{\ifanon\textit{[link withheld for double-anonymous review]}\else\href{#1}{#2}\fi}

\SetKwInOut{Input}{Input}\SetKwInOut{Output}{Output}\SetKwInOut{Data}{Data}
\SetKwProg{Fn}{Function}{:}{}
\SetKwComment{Comment}{$\triangleright$\ }{}

\begin{document}

% ---- title and authors ---------------------------------------------------------------------------
% TODO(title): descriptive, no "new"/"novel", no maths symbols, Title Case (checklist section 2). Working title placeholder only.
\title{[AUTHOR: title, to be set when the method and claim are frozen]}

\ifanon
  \author{Anonymous authors for double-anonymous review%
  \thanks{Manuscript submitted to the IEEE Transactions on Neural Networks and Learning Systems.}}
\else
  % Camera-ready / non-anonymous: first footnote carries dates, corresponding author, FULL financial support, affiliations, e-mails.
  % TODO: add \IEEEmembership{Member, IEEE} after a name only if that author is an IEEE member (affiliate grades are not listed).
  \author{Shlok~Kumar~Goenka and [AUTHOR: supervising professor, full name]%
  \thanks{[AUTHOR: funding statement or ``This research received no external funding.'']}%
  \thanks{(Corresponding author: [AUTHOR: corresponding author].)}%
  \thanks{The authors are with [AUTHOR: school or department], Vellore Institute of Technology, Vellore 632014, Tamil Nadu, India (e-mail: [AUTHOR: e-mails]).}%
  \thanks{[AUTHOR: ORCID of each author is entered in the Author Portal; it is not printed in the PDF unless the template asks.]}}
\fi
% IEEEtran journal mode: running head (page numbers sit in the head). Anonymous build keeps the head free of author names.
\markboth{IEEE Transactions on Neural Networks and Learning Systems}{Anonymous: [AUTHOR: short title]}

\maketitle

% ---- abstract and index terms --------------------------------------------------------------------
% TODO(abstract): ONE paragraph, 150-250 words, no abbreviations, footnotes, references or equations; write last.
\begin{abstract}
% abstract.tex -- single-study version, 2026-10-31 (Shlok approved: one study, abstract without numbers; PROFESSOR-SUGGESTIONS-ANALYSIS-2026-10-31.md).
% No number, interval, p-value, sample size or parameter count; every number stays in Results. One paragraph, no abbreviation, symbol,
% equation or citation. Deviation from the A6.12 sentence templates (numbers moved to Results) is listed in the Deviations section.
% Claims and their sources (all in Results, with run comments there):
% run: runs/P2-LRF-L1/a6_analysis.json:family_L.L1.label (= confirmed); family_L.L-LEN.label (= rejected); family_L.L-NI.{non_inferior (= true), margin_note}
% run: runs/P2-LRF-DN/dn_analysis.json:cells.{noise0.1,jpeg10}.arms.{DN+NONE,DN+LRF}.vs_LRF (denoise before scoring: above LRF under noise; DN+LRF below LRF under JPEG; DN+NONE under JPEG interval contains 0)
% run: runs/P2-LRF-ABL-SCORE/ablation_analysis.json:rows[id=A-LIN] (linear logit: noise interval contains 0; at or above LRF in the other three conditions)
Image-text verification, the choice of the true caption of an image from two candidates, becomes harder when the image is damaged.
We study this problem inside one frozen pretrained unified discrete diffusion model, UniDisc, which gives every observed image token a
probability given a caption, with no external encoder.
For each token, the log-probability ratio between the true and a wrong caption is evidence, and the standard rule adds the evidence of
all tokens with equal weight.
We learn a per-token reliability weighting of this evidence from features computed without the caption, fit it on design images, and
freeze it.
A preregistered one-sided exact paired test on fresh Common Objects in Context validation images under Gaussian noise at one level
confirmed that the frozen weighting raised verification accuracy, including when the wrong caption was matched to the true one in length.
On clean images the weighting was non-inferior within a preregistered margin; the margin was fixed after an exploratory clean-image
interval had been seen.
Two exploratory comparisons qualify this result.
A classical pixel denoiser applied before scoring reached higher accuracy than the weighting under Gaussian noise, and combined with the
weighting it reached lower accuracy under Joint Photographic Experts Group compression.
A linear logit on the same features reached similar accuracy under Gaussian noise and higher accuracy under compression.
An exploratory sweep on the same images found the gain at milder and stronger Gaussian noise as well.
The confirmatory evidence covers one model, one dataset, and one noise level; it does not show that the learned weights track token
reliability, or that the gain holds for other models or corruptions.

\end{abstract}

% TODO(index terms): 4-5 terms, alphabetical, first word capitalised, from the IEEE Thesaurus where possible.
\begin{IEEEkeywords}
[AUTHOR: index term 1], [AUTHOR: index term 2], [AUTHOR: index term 3], [AUTHOR: index term 4].
\end{IEEEkeywords}

% ---- body ----------------------------------------------------------------------------------------
\section{Introduction}
\label{sec:intro}
% TODO(introduction): positions the work in the literature, states the question and why it matters, lists contributions,
% gives a section roadmap. Cite TNNLS papers only where each genuinely supports its sentence (Shlok, memory: cite-target-journal-papers).
\todo{TODO: introduction}
      % I   Introduction
% related.tex -- Section II (writer W1, 2026-10-30). Binding contract: paper2-tnnls/SPEC.md v2.2, section 3 row II.
% Every citation is a verified key of paper2-tnnls/references.bib, used for at most the sentence its "supports" line in
% research/paper2/tnnls/CITATION-AUDIT-2026-10-30.md allows (sections 1-2). The eight TNNLS papers: moser2025diffusion, yang2026malfm,
% li2026tcpfmc, chen2025uncertainty, zhang2024informative, jiang2025resisting, wang2025conformal, qiu2025counterfactual.
% Study 2's novelty sentence is written for Study 2 alone and scoped to the logged search (NOVELTY-AUDIT-LRF-2026-10-29.md sections 1,
% 4, 5, 7); the audit's matched-filter sentence and its section-6 sentence are NOT used (F13). Lumina-DiMOO is not called masked
% discrete diffusion (unverified). Muddit is named once with its published title (citation audit, open item 3).

\section{Related Work}\label{sec:related}

\paragraph*{Generative and likelihood scoring}
Diffusion models are trained by maximising a variational lower bound on the data log-likelihood, as surveyed for continuous diffusion
in image super-resolution~\cite{moser2025diffusion}, so a trained model can be read as a scorer of its own inputs.
Text-to-image diffusion models have been used as zero-shot classifiers whose weights over diffusion timesteps are learned from held-out
labelled data~\cite{clark2023zeroshot}, including with few labels~\cite{jeong2025diffusion}.
For image-text matching, the conditional error has been normalised by the unconditional one~\cite{krojer2023diffusionitm}, and VQAScore
scores alignment by the probability of ``Yes'' that an image-to-text model assigns to a question about the pair~\cite{lin2024vqascore}.
In text, masked language models score a token with the token masked~\cite{salazar2020mlm,li2022uchecker}.
The evidence used in this paper belongs to this family: a per-token log-likelihood ratio of the observed image codes under two captions,
computed by the masked denoiser of a unified model.

\paragraph*{Masked discrete diffusion and unified models}
Absorbing-state discrete diffusion~\cite{austin2021d3pm,lou2024sedd} was simplified into masked diffusion language
models~\cite{sahoo2024mdlm} and scaled~\cite{nie2025llada}; the denoiser of such a model represents the conditional distribution of a
clean token given the unmasked context~\cite{ou2025radd}.
MaskGIT generates image tokens by masked prediction~\cite{chang2022maskgit}, and unified text-image models place several modalities in one
sequence~\cite{swerdlow2025unidisc,yang2025mmada,li2026omnidiffusion}, as does Muddit, published as ``Beyond text-to-image: Liberating
generation with a unified discrete diffusion model''~\cite{shi2025muddit}.
LaViDa adds a vision encoder to a diffusion language model~\cite{li2025lavida}, and Lumina-DiMOO is a further omni diffusion language
model for multimodal generation and understanding~\cite{xin2025luminadimoo}.
In this journal, discrete-diffusion image captioning has been studied in a setting where noise in the diffusion process weakens
sentence information, and a feature-mask and gated-fusion design improves image-text alignment~\cite{yang2026malfm}.
Many methods remask tokens that the sampler itself produced, by confidence, learned quality scores, or critics, including
cross-modal confidence in multimodal masked diffusion models~\cite{wang2025remdm,zhai2026core,kim2025prism,schiff2026proseco,lezama2022tokencritic,ouyang2026freecorrection,le2026co2jump}.
Work on damaged \emph{input} tokens is scarcer: Trust but Verify detects corrupted context tokens in a finetuned text-only masked
diffusion model and refills them~\cite{jhingran2026trustverify}, and dropout decoding masks visual tokens that are uncertain under a
logit-lens projection~\cite{fang2025dropout}.
None of them reports a verification decision under image corruption in a unified masked diffusion model.

\paragraph*{Reliability-aware and dynamic fusion}
When evidence comes from sources of unequal quality, the classical combination rule weights each source by its
reliability~\cite{chair1986optimal}, and classifier combination studies the rules for doing so~\cite{kittler1998combining}.
In multimodal learning the weights change per input: dynamic, uncertainty-driven weighting of modalities has been analysed for
low-quality data, with theory on when it is preferable to static weighting~\cite{zhang2023provable}. A recent survey organises fusion on low-quality data into noisy, incomplete, imbalanced, and
quality-varying settings~\cite{zhang2026lowquality}.
Trusted multi-view classification fuses per-view evidence dynamically to handle unreliable views~\cite{han2023trusted}, and evidential
methods attach an estimated importance or reliability to each modality before combining~\cite{shao2024ddef,huang2025evidential}.
Papers in this journal take the same view: a per-modality confidence score (an energy score) lets the fusion rely less on low-quality
modalities~\cite{li2026tcpfmc}; evidence theory quantifies the uncertainty of each prediction and fuses evidence across propagation
depths~\cite{chen2025uncertainty}; evidential deep learning suppresses unreliable modality-wise pseudo-labels in audio-visual video
parsing~\cite{jiang2025resisting}; and counterfactual samples in a shared multimodal latent space identify the input features that
contribute to high predictive uncertainty~\cite{qiu2025counterfactual}.
At a finer grain, attention-based multiple-instance pooling learns instance weights from instance features under a bag-level
loss~\cite{ilse2018attention}, and a token weight has been derived from a model's logits with and without the image, as a training-time
reweighting~\cite{xiao2024seeing}.
For Study~2, the closest precedents are thus learned timestep weights~\cite{clark2023zeroshot,jeong2025diffusion}, instance weights
under a bag-level loss~\cite{ilse2018attention}, and reliability weights of modalities or views~\cite{zhang2023provable,han2023trusted}.
Token-level confidences of a captioning vision-language model, aggregated algebraically or by a learned estimator, have also been used to
judge whether a caption is correct for its image~\cite{petryk2024token}; that work scores the caption's tokens given the image and does
not study image corruption.
To our knowledge, no earlier work learns per-token, caption-withheld weights for the log-likelihood-ratio evidence of a unified
masked diffusion model and tests them on fresh images under image corruption.
This statement is limited to a keyword search of arXiv and of title-level records of IEEE-published journals and proceedings, logged in
October 2026; the search did not cover full texts, the pages of IEEE Xplore itself, or the submission lists of open-review venues.

\paragraph*{Conformal risk control}
Conformal prediction gives distribution-free guarantees under exchangeability~\cite{vovk2005algorithmic,angelopoulos2023gentle};
conformal $p$-values give false-discovery control for outlier detection~\cite{bates2023outliers}, and conformal risk control bounds the
expected value of a monotone loss~\cite{angelopoulos2024crc}.
Conformal loss-controlling prediction, published in this journal, extends conformal prediction from miscoverage to a general loss for
each test object, with a finite-sample guarantee under exchangeability~\cite{wang2025conformal}; Study~1 uses the expected-loss
statement of conformal risk control instead.
Conformal risk control has been applied, with masking-derived CLIPScore distributions and a threshold set on foil-labelled data,
to flag erroneous caption words given an image~\cite{gomes2025crc}.

\paragraph*{Corruption benchmarks and preregistration}
Common-corruption benchmarks test models under non-adversarial image corruptions~\cite{hendrycks2019benchmarking}, and a
corruption benchmark for object detection released an implementation of such corruptions at graded
severities~\cite{michaelis2019imagecorruptions}, which Study~2 uses for its other noise types.
Preregistration records the questions and the analysis plan before the outcomes are observed, which separates prediction from
postdiction~\cite{nosek2018preregistration}; the form used here is defined in \cref{sec:setting-tests}.
           % II  Related Work
\section{Preliminaries}
\label{sec:prelim}
% TODO(preliminaries): notation table or paragraph; problem statement; assumptions as \begin{assumption} environments.
% Every symbol is defined at first use; every equation is numbered, referred to with \cref, and motivated in words.
\todo{TODO: notation and problem statement}
     % III Preliminaries
\section{Method}
\label{sec:method}
% TODO(method): the method is being designed (2026-10-29); nothing is written here on purpose.
% Algorithms: algorithm2e (ruled, vlined, linesnumbered) with Input, Output and Data blocks, built up across the paper so that the
% final algorithm calls the earlier ones. Keep the main text to the algorithm that matters; others go to the supplement.
\todo{TODO: method}
            % IV  Method
\section{Theory}
\label{sec:theory}
% TODO(theory): statements (theorem/proposition/lemma environments, defined in main.tex) with their assumptions and what is proved
% versus what is only measured. Long proofs go to the appendix (counts toward 15 pages) or to supplement/supplement.tex (does not).
\todo{TODO: theory}
            % V   Theory
\section{Experimental Setup}
\label{sec:setup}
% TODO(setup): models, datasets, corruptions, baselines, metrics, protocol (preregistration reference, anonymised link in review),
% seeds, hardware, hyperparameter budgets. Each empirical statement carries a run ID in a LaTeX comment: % run: <exp-id>
\todo{TODO: experimental setup}
             % VI  Experimental Setup
\section{Results}
\label{sec:results}
% TODO(results): main comparison, ablations, robustness, cost/complexity. No claim stronger than its confidence interval; negative
% results and losses to a baseline stay in. Every number traces to runs/<exp-id>/  (% run: <exp-id> next to the sentence).
% Figures: vector PDF, one column 3.5 in or two columns 7.16 in, fonts 8--10 pt, readable in greyscale; captions only in LaTeX.
\todo{TODO: results}
           % VII Results
% discussion.tex -- Section VI (writer W1, 2026-10-30). Binding contract: paper2-tnnls/SPEC.md v2.2, section 3 row VI, sections 7, 8.
% VI-A opens with the A4.2.3 NEG "6.1 opens" sentence; the Tier-3 paragraph and the family H paragraph come after it (A4.1.10 placement
% 2; A5.9 placement 2: own paragraph headed with the prefix, no sentence about P). VI-B is Study 2 alone; the alternatives are the red
% team's items 1-5 (REDTEAM-SPEC "Cheapest alternative explanations"), worded as open questions; item 6 is withdrawn (V2-1) and not used.
% VI-C follows SPEC section 7 item by item (no sentence about the P discordance, V2-1); VI-D follows SPEC section 8.
% No sentence links the two studies; no comparison word; no theory word for a result; every number has a run comment.

\section{Discussion}\label{sec:discussion}

\subsection{Study 1}\label{sec:disc-study1}
% run: runs/P2-GATE-M4/analysis.json:outcome.label (= negative); primary.per_corruption.*.ci90
The preregistered test excludes an improvement or a loss of more than five percentage points in image-direction verification accuracy
from gating under JPEG quality~10, Gaussian noise with $\sigma=0.1$ and a 128-pixel splice, for this model and these severities; it does
not show that gating has no effect.
It concerns this gate at its preregistered level $\alpha=0.05$, one severity per corruption, one model and one data set, and says
nothing about other levels, stronger corruptions, models or tasks, none of which was tested.

% run: runs/P2-GATE-M4/tier3_fusion.json:cells.*.contrasts.* (exploratory; Section IV-C)
The exploratory Tier-3 comparison of the hard gate with five combiners of the same evidence is reported in \cref{sec:study1-results} and
supplement S-II; it describes the combiners and is not an account of the primary result.

% run: runs/P2-GATE-H3/family.json:{n_headroom_items (= 50), n_rule.h_rule (= H-NOGO), n_rule.G_max_headroom (= null), go (= false)}
\emph{In a separately labelled family added after the primary result was known (family H: SugarCrepe hard
negatives~\cite{hsieh2023sugarcrepe} on 50 COCO images not used elsewhere; Holm over the same three corruptions),} the pre-written pilot
rule found no corruption with decision headroom on 50 pilot images, so the family's test was not run, and nothing is stated about the
gate on hard negatives (\cref{sec:study1-H}; supplement S-III).

\subsection{Study 2}\label{sec:disc-study2}
% run: runs/P2-LRF-L1/a6_analysis.json:family_L.L1.label (= confirmed); branch.branch (= A)
L1 shows that one frozen weighting, hash-pinned before any test image was chosen, raised image-direction verification accuracy under
Gaussian noise with $\sigma=0.1$ on 450 COCO images that played no part in its design or fit, passing an exact paired test at the
preregistered level and holding with length-matched negative captions.
It does not show that the weighting works on a corruption it was not trained on, that its weights follow the reliability of a token,
or that the effect carries to another model, data set or severity.
% run: runs/P2-LRF-LG/a6_analysis.json:family_LG.cells.impulse_noise3.ci95 (= [0.0, 0.1]); runs/P2-LRF-L1/a6_analysis.json:family_L.L4.p_holm (= 0.0486)
Five inexpensive alternative readings remain open, each with its evidence in \cref{sec:limitations}: (i) a reweighting that helps
whatever the condition; (ii) a supervised fit to the training corruption; (iii) a single noise-sensitive feature, such as the
codebook-neighbour distance or the one-pass observed-position score, which \cref{fig:lrf_weights} cannot separate and the exploratory
ablation of \cref{sec:study2-ablation} bears on only as an exploratory result; (iv) a limited reach to other noise types, since the
impulse-noise interval starts at 0.0; and (v) an unstable splice decision, whose Holm-adjusted $p$ of 0.049 lies 0.001 below the level.
A second backbone, a corruption held out of training and a preregistered comparison with simpler weighting rules on fresh images are the
tests that would settle these questions; none was run.

\subsection{Limitations}\label{sec:limitations}
\emph{Study 1.} One model and one severity per corruption were tested.

% run: runs/P2-LRF-L1/a6_analysis.json:family_L.L4.p_holm (= 0.0486); runs/P2-LRF-LG/a6_analysis.json:family_LG.cells.impulse_noise3.p_two (= 0.058)
% run: runs/P2-LRF-L1.try1/selfcheck.json:image_tokens_frac_equal (= 0.98828125); runs/P2-LRF-L1.try2/selfcheck.json:image_tokens_frac_equal (= 0.9892578125)
% run: runs/P2-LRF-BASELINES-PARITY-LG/baselines.json:pixel_parity.min (= 0.986328125)
\emph{Study 2.} One model, one data set and one severity per corruption were tested.
LRF was trained on the tested corruption (Gaussian noise with $\sigma=0.1$ was a training condition), so L1 tests new images, not a new
corruption.
No preregistered test separates corruption-specific weighting from a generic gain, since the clean-image and splice estimates were also
positive; the title's ``under Gaussian noise'' names the tested condition, not a specificity.
The Holm-adjusted $p$ of L4 is 0.049, and the two-sided $p$ of the impulse-noise cell is 0.058.
The clean-image margin of L-NI was fixed after an exploratory clean interval had been seen.
Whether COCO images were in the pretraining data of UniDisc, its tokeniser, CLIP or SigLIP is unknown.
Each image has one corruption realisation and one mask seed, so the intervals cover item sampling only.
A fresh tokenisation of the same pixels in another process need not reproduce the stored codes: it agreed on 98.8\,\% and 98.9\,\% of
image tokens in the two stopped self-checks of card row A-8, and on at least 98.6\,\% in the pixel-parity runs.
The reference encoders almost solve this matching task (\cref{sec:study2-reference}); this paper makes no claim of a verifier that ranks
with them.
No ablation was preregistered; \cref{sec:study2-ablation} is exploratory.

\emph{Both studies.} All results are for random-caption negatives. The preregistration is self-hosted in the sense of
\cref{sec:setting-tests}, and the audits were agent-run (see the Declaration).

\subsection{Deviations From the Plans}\label{sec:deviations}
{\small
\noindent\emph{Study 1:}
(1) The Tier-3 comparators, parameters and fallback seed were fixed in an addendum committed before the main run's headroom decision;
their intervals carry no $p$-value and no multiplicity adjustment.
(2) Outcome-label wording was fixed before the primary result, with a dated record of the null wording and of the reading adopted where
card and plan differed on a rejection against the gate.
(3) The time-input statement rests on a code reading at a pinned commit and a key check; its table moved to the supplement.
(4) Family~B ran before the re-infilling family under the plan's fixed time inequality.
(5) Family~H has its own disclosure and amendment items; one decision row was committed after its pilot driver started and before any
pilot row existed.
(6) The effective-$n$ record of the re-infilling figure carries a reason text written at analysis time.
(7) The data steward did not collect the image ids of negatives and mismatches separately (no analysis uses them), and its check left
some run files outside the excluded classes unopened.
(8) The exact Clopper--Pearson bounds of \cref{tab:familyP} were added at writing (\nolinkurl{runs/P2-PAPER-BOUNDS}, descriptive).
\emph{Study 2:}
(1) The L1 stage stopped twice at its tokenisation self-check before any item; card row A-8 moved the code pin, and row A-9 changed one
test only.
(2) The scorer's pin file was written while L1 ran, with every value copied from records committed before L1.
(3) The environment lock files of the design, re-score, L1 and other-noise stages are empty; the environment was recorded after the runs,
and the corruption package's commit was verified from it.
(4) Item arrays are untracked with committed SHA-256 manifests; hashes appear in CRLF and LF forms of the same content.
(5) The pixel-parity files were moved unchanged from a wrong output directory with no prior hash (audit-time hashes are recorded); they
supersede the field ``pixel\_parity: not run'' of the main reference-encoder file.
(6) The reference-encoder download was larger than the size quoted to the first author.
(7) What had been seen before the re-freeze is stated in \cref{sec:study2-provenance}.
(8) ``No output was read before its run was complete'' rests on the operations log.
(9) The exploratory ablation was added at writing.
Supplement S-I gives each item in full with its record.\par}
        % VIII Discussion and Limitations
% conclusion.tex -- Section VI (single-study version, 2026-10-31; two-study text at git tag paper2-two-study-2026-10-31).
% One paragraph for the one study; no link to any other study; no rescue; no ablation, comparison or reference-encoder content
% (SPEC section 6). Numbers repeat the binding sentences, each with its run comment.

\section{Conclusion}\label{sec:conclusion}
% run: runs/P2-LRF-L1/a6_analysis.json:family_L.L1.{acc_b (= 0.7333), acc_a (= 0.8022), net (= 0.0689), ci95 (= [0.0289, 0.1089]), p_one (= 0.000508)}; family_L.L-LEN.label (= rejected)
% run: runs/P2-LRF-L1/a6_analysis.json:family_L.L-NI.{ci90 (= [0.0089, 0.0533]), margin (= 0.03), non_inferior (= true)}; runs/P2-LRF-PLAN/n_rule.json:n_L (= 450)
We tested one frozen, learned reliability weighting of the per-token evidence of a pretrained unified masked diffusion model on a fresh
preregistered fold.
Under Gaussian noise with $\sigma=0.1$ it raised image-direction verification accuracy from 0.733 without reliability weighting to 0.802
(net 0.069, 95\% interval 0.029 to 0.109; exact one-sided McNemar $p=0.0005$), including with length-matched negatives, and on clean
images it lost no more than three percentage points (0.03 in accuracy; 90\% interval of the net change 0.009 to 0.053); the margin was
fixed after an exploratory clean interval had been seen.
% run: runs/P2-LRF-F30K/f30k_analysis.json:F1.{acc_none (= 0.6927), acc_lrf (= 0.7924), net (= 0.0998), ci95 (= [0.0722, 0.1284]), p_one (= 6.44e-12), label (= confirmed)}; F_LEN.label (= confirmed); F_NI.{lower90 (= 0.0092), margin (= -0.03), non_inferior (= true)}; n_scored (= 872)
A second preregistered test, on 872 Flickr30K test images with the weighting unchanged, confirmed the gain under the same noise (0.693 to
0.792; net 0.100, 95\% interval 0.072 to 0.128; exact one-sided McNemar $p=6.4\times10^{-12}$), also with length-matched wrong captions, and
the weighting was non-inferior on clean images.
This holds for one model, on 450 COCO images and 872 Flickr30K images not used to design or fit the method, at one severity; the weighting
was trained on the tested corruption, and no preregistered test shows that its gain is specific to that corruption.
        % IX  Conclusion

% ---- appendix (counts toward the 15-page limit; put long material in supplement/supplement.tex instead) ----
\appendices
% appendix.tex -- Appendix A: proof of prop:snrcorr only (SPEC v2.2; THEORY-ADDENDUM r1 R1.8 appendix block, in the notation of Section IV).
% Prop. 3, Lemma 2 and Cor. 3.1 are not here: supplement S-II only (D0090).
\section{Proof of \Cref{prop:snrcorr}}\label{app:snrcorr}
Write $m_0=(1-\rho)m$ and $m_1=\rho m$ for the numbers of clean and corrupted tokens, $k_0=1-\phi_0$ and $k_1=1-\phi_1$ for the keep
probabilities, and $z_j\in\{0,1\}$ for the label of token $j$, so that $\E\Delta_j$ is $\mu$ or $\mu_c$ and $\Var\Delta_j$ is $\sigma^2$ or
$\sigma_c^2$.
The exchangeable block covariance of assumption (a) is positive semi-definite if and only if $\gamma_0\ge-1/(m_0-1)$,
$\gamma_1\ge-1/(m_1-1)$, $\gamma_0,\gamma_1\le1$ and $\gamma_{01}^2\,m_0m_1\le\{1+(m_0-1)\gamma_0\}\{1+(m_1-1)\gamma_1\}$.

\paragraph*{Variance}
By assumption (b), $I_j$ is independent of $\Delta_j$, so
$\Var(I_j\Delta_j)=k(\sigma_z^2+\mu_z^2)-k^2\mu_z^2=k\sigma_z^2+k(1-k)\mu_z^2$ with $k=k_{z_j}$, which are the terms of \cref{eq:var}.
For $j\neq k$, again by (b),
$\mathrm{Cov}(I_j\Delta_j,I_k\Delta_k)=\E I_j\,\E I_k\,\E[\Delta_j\Delta_k]-\E I_j\,\E I_k\,\E\Delta_j\,\E\Delta_k=k_{z_j}k_{z_k}\,\mathrm{Cov}(\Delta_j,\Delta_k)$.
There are $m_0(m_0-1)$ ordered clean pairs with covariance $\gamma_0\sigma^2$, $m_1(m_1-1)$ ordered corrupted pairs with covariance
$\gamma_1\sigma_c^2$ and $2m_0m_1$ ordered mixed pairs with covariance $\gamma_{01}\sigma\sigma_c$, which gives \cref{eq:varcorr}.

\paragraph*{Condition}
Let $V_g$ be \cref{eq:varcorr} at $(\phi_0,\phi_1)$ and $V_0$ its value at $\phi_0=\phi_1=0$.
With $\mu_c=0$, $\E S_{\rm G}=m_0k_0\mu$ with the gate and $m_0\mu$ without it, so, both being positive, the gate raises the ratio if and
only if $k_0^2V_0>V_g$.
Collect $k_0^2V_0-V_g$ by type of term.
The clean terms are
$k_0^2m_0\sigma^2\{1+(m_0-1)\gamma_0\}-m_0k_0(\sigma^2+\phi_0\mu^2)-m_0(m_0-1)k_0^2\gamma_0\sigma^2=-m_0k_0\phi_0(\sigma^2+\mu^2)$,
in which the $\gamma_0$ terms cancel.
The corrupted terms are $m_1\sigma_c^2\{(k_0^2-k_1)+(m_1-1)\gamma_1(k_0^2-k_1^2)\}$, with $k_0^2-k_1=\phi_1-2\phi_0+\phi_0^2$ and
$k_0^2-k_1^2=(\phi_1-\phi_0)(2-\phi_0-\phi_1)$.
The mixed terms are $2m_0m_1\gamma_{01}\sigma\sigma_c\,k_0(k_0-k_1)$, with $k_0-k_1=\phi_1-\phi_0$.
Dividing the sum by $m$ gives \cref{eq:snrcorr}.

\paragraph*{Consequences}
(i) is read off \cref{eq:snrcorr}.
(ii) At $\phi_1=\phi_0\in(0,1)$ every correlation term vanishes and the left side is $\rho\sigma_c^2(\phi_0^2-\phi_0)<0$, below the positive
right side.
For $\phi_1\le\phi_0$, let $G=(\rho m-1)\gamma_1$, so that $G\ge-1$ by positive semi-definiteness.
The bracket $(\phi_1-2\phi_0+\phi_0^2)+G(\phi_1-\phi_0)(2-\phi_0-\phi_1)$ is linear in $G$; it equals $\phi_1-2\phi_0+\phi_0^2\le-\phi_0(1-\phi_0)$
at $G=0$, equals $-\phi_1(1-\phi_1)\le0$ at $G=-1$, and does not increase for $G\ge0$; hence it is at most $0$.
If $\gamma_{01}\ge0$ the mixed term is at most $0$ as well, while the right side is positive for $\phi_0\in(0,1)$.
(iii) The derivatives of the left side with respect to $\gamma_1$ and $\gamma_{01}$ are
$\rho\sigma_c^2(\rho m-1)(\phi_1-\phi_0)(2-\phi_0-\phi_1)$ and $2\rho(1-\rho)m\sigma\sigma_c(1-\phi_0)(\phi_1-\phi_0)$, both positive when
$\phi_1>\phi_0$ and $\rho m>1$, and the right side does not depend on $\gamma_1$ or $\gamma_{01}$.
For fixed $\phi_1>\phi_0$ the $\gamma_1$ term grows without bound with $(\rho m-1)\gamma_1$, which gives the limit.
With $\gamma_{01}=0$, $x=\phi_1-\phi_0$ and $\chi=\rho\sigma_c^2/\{(1-\rho)(\sigma^2+\mu^2)\}$, \cref{eq:snrcorr} reads
$Gx^2-x\{1+2G(1-\phi_0)\}+\phi_0(1-\phi_0)(1+1/\chi)<0$, so for $G>0$ the gate raises the ratio if and only if $\phi_1$ exceeds
$b(\phi_0;G,\chi)=\phi_0+\big[1+2G(1-\phi_0)-\sqrt{\{1+2G(1-\phi_0)\}^2-4G\phi_0(1-\phi_0)(1+1/\chi)}\big]/(2G)$ (and $\phi_1\le1$), the curve of
\cref{fig:frontier}; as $G\to0$ it tends to the boundary $2\phi_0-\phi_0^2+\phi_0(1-\phi_0)/\chi$ of \cref{eq:snr}.
(iv) For $\gamma_{01}<0$ and $\phi_1<\phi_0$ the mixed term is positive and grows with $m$, so it can exceed the remaining terms.
% run: runs/P2-THEORY-CHECK/output.txt (check D4: min eigenvalue 0.9500; SNR none 0.4892, gate (0.3, 0) 0.5080)
An instance with a positive definite covariance is $m_0=900$, $m_1=100$, $\gamma_0=\gamma_1=0.05$, $\gamma_{01}=-0.045$, $\mu=0.1$, $\mu_c=0$,
$\sigma^2=\sigma_c^2=1$: the smallest eigenvalue of the correlation matrix is 0.9500, and the ratio rises from 0.4892 without a gate to
0.5080 with the gate $(\phi_0,\phi_1)=(0.3,0)$. The instance illustrates the algebra and is not an estimate from the model.
(v) $m$ enters \cref{eq:snrcorr} through $\rho m-1$ and $m$ whenever $\gamma_1\neq0$ or $\gamma_{01}\neq0$. \qed


% ---- Declaration of Generative AI Use ------------------------------------------------------------
% Shlok's rule: AI use is mentioned ONLY in this section, and there is NO Acknowledgment section.
% IEEE's rule (CIS Information for Authors K; Author Center submission policies) puts the disclosure in Acknowledgments and asks
% for a citation to the AI system in the affected sections. The conflict is recorded as C2 in the checklist; not resolved here.
% declaration.tex -- unnumbered back matter before the References (writer W1, 2026-10-30).
% (1) Reproducibility and Data Availability: SPEC v2.2 section 9 (claims nothing beyond the listed facts; F16).
% (2) Declaration of Generative AI Use: SPEC rule 6, D0093 (the ONLY place AI use is mentioned; no Acknowledgment section). Written so
%     that it can be moved unchanged into an Acknowledgment if the editor asks (checklist C2): it names each system and version, the
%     roles, the sections affected and the level of use (IEEE Author Center, AI-generated content policy). Codex model name from
%     research/paper2/tnnls/codex/SESSION.txt. Facts only the authors can confirm are marked \todo{AUTHOR: ...}.

\section*{Reproducibility and Data Availability}
The code, the preregistration documents (cards, analysis plans and their addenda), the hash pins, the stored predictions and the analysis
outputs are kept in one version-controlled repository
(\anonlink{https://github.com/shlokDS16}{repository}\ifanon\else\ \todo{AUTHOR: exact repository URL and archive DOI}\fi);
% TODO(AUTHOR): replace the \anonlink target with the exact repository URL (and an archive DOI) before the camera-ready build.
for review, an anonymised mirror is provided that preserves the commit times and hashes on which the definition of ``preregistered'' in
\cref{sec:setting-tests} rests.
The per-item arrays are not tracked; their SHA-256 manifests are committed.
The software environment was recorded after the runs.
The images and captions are the public COCO val2017 split~\cite{lin2014coco}, and the model is the released UniDisc
checkpoint~\cite{swerdlow2025unidisc} at the snapshot pinned in supplement S-I.

\section*{Declaration of Generative AI Use}
The authors used generative AI systems in this work.
The systems were Claude Opus~5.5, Claude Sonnet~5.5 and Claude Haiku~5.5 (Anthropic), run as agents through Claude Code, and Codex
(OpenAI; model recorded in the session log as \texttt{gpt-6-luna}), used to review plans and drafts.
\todo{AUTHOR: confirm that this list of systems and versions is complete for every stage of the work, and give the Claude Code and Codex
client versions if the editor requires them.}
They were used in four roles.
(i) Drafting and editing the text of every section, the abstract, the figure and table captions and the supplement.
(ii) Writing the experiment code, the analysis scripts and the figure scripts, and producing the figures and tables from the committed
run files.
(iii) Literature search and the verification of every reference against Crossref, arXiv, OpenAlex or the publisher's record.
(iv) The scripted and agent-run audits named in the paper: the red team (adversarial review of the plans, results and text), the
gate-auditor (commit order, hash checks, code audits and reconciliation of every number with the run files), the data-steward (fold
construction and leak checks) and the statistician (analysis plans, statistical checks and the ruling on the abstract).
The sections affected are therefore all sections, the supplement, the figures and the tables.
The level of use was substantial: the systems wrote first drafts of the text and code and ran the audits, and the authors directed the
work, set the research questions and approved the preregistration documents and the decisions recorded in them.
\todo{AUTHOR: confirm the level-of-use sentence, in particular which decisions the authors made and approved themselves.}
No AI system generated data: every reported number comes from runs of the stated model on the stated images.
The audits listed above were run by AI agents inside the authors' own pipeline, not by outside reviewers.
The authors reviewed all content and take full responsibility for it.


% ---- references ----------------------------------------------------------------------------------
\ifplaceholder\nocite{*}\fi
\bibliographystyle{IEEEtran}
\bibliography{references-ieee}      % generated from references.bib by scripts/prepare_bib.py (DOI/arXiv -> url, links print blue)

% ---- biographies and photos: final (non-anonymous) stage only -----------------------------------
\ifanon\else% Biographies and photographs: final (non-anonymous) stage only. Not included in the anonymous build (see main.tex).
% IEEE: up to three paragraphs; bold full name matching the byline; education with years; positions; head-and-shoulders photo;
% QR codes are not accepted instead. Photo files: first five letters of the surname (JPEG allowed for photos), >300 dpi.
\begin{IEEEbiographynophoto}{Shlok Kumar Goenka}
[AUTHOR: biography, education with years, positions, current affiliation; photo via the IEEEbiography environment.]
\end{IEEEbiographynophoto}
\begin{IEEEbiographynophoto}{[AUTHOR: supervising professor, full name]}
[AUTHOR: biography.]
\end{IEEEbiographynophoto}
\fi

\end{document}
